\documentclass[leqno, a4paper, 12pt]{jsarticle}
\usepackage{amssymb}
\usepackage{amsthm}
\usepackage{amsmath}
\usepackage{comment}
\usepackage{url}

\usepackage{enumitem}
%\usepackage{showkeys}


%定理環境 thmはあまり使わずpropをなるべく使う。
%主張は基本的に証明の中で使い、そのため番号を付けない。
%注意環境を付けてみた。
\newtheorem{thm}{定理}
\newtheorem{prop}[thm]{命題}
\newtheorem{cor}[thm]{系}
\newtheorem{lem}[thm]{補題}
\newtheorem{df}[thm]{定義}
\newtheorem{exam}[thm]{例}
\newtheorem{fact}[thm]{事実}
\newtheorem*{claim}{主張}
\newtheorem*{cau}{注意}
\newtheorem*{rmk}{注意}
\renewcommand{\proofname}{証明}
%矢印たち。
%\toは\rightarrowと同じ。\getsは\leftarrowと同じ。同値は\iff
\newcommand{\To}{\Longrightarrow}
\newcommand{\Gets}{\LongLeftarrow}
%黒板太字
\newcommand{\nn}{\mathbb{N}}
\newcommand{\zz}{\mathbb{Z}}
\newcommand{\qq}{\mathbb{Q}}
\newcommand{\rr}{\mathbb{R}}
\newcommand{\cc}{\mathbb{C}}
%無限大
\newcommand{\mgn}{\infty}
%差集合は\setminusで統一する。等号の否定は\neq
%真部分集合は\subsetneqq　偏微分は\partial
%ディスプレイスタイル。\dfracでディスプレイ分数
\newcommand{\dpst}{\displaystyle}

\title{論文メモ
\large{\\--- Urysohnの論文 ---}
}
\author{ナンブ　キトラ}
\date{\today}
\begin{document}
\maketitle

本棚を整理したいので，
溜まっている論文のメモを書いて未来への手紙とすることにする．

この論文の束はウリゾーンの論文を集めたものだと思う．
全部ドイツ語で書かれているのでマジで読めない．
1924年のウリゾーンの死後にアレキサンドロフが代理出版していたりするので，著者がアレキサンドロフになっていたりすることがある．
グーグル翻訳のカメラ機能を使ってなんとか読もうと努力した．

論文\cite{MR1512214}
は読めなかった．なんかコンパクト空間の可算性と非可算性を研究しているような気がする．
コンパクト空間の閉集合が$G_{\delta}$
ならば第一可算公理を満たすと言うようなことを証明しているような気がする．

論文
\cite{MR1512259}
は読めなかった．
何やら空間が$\rr^{n}$
に同相?位相的埋め込み?であるための十分条件を与えているような気がする．


論文
\cite{MR1512213}
は読めなかった．
完全集積点，つまりその点の任意の近傍の濃度は全体空間の濃度に一致するような点のことだが，これとコンパクト性の関係性を見ているような気がする．昔の人がいうコンパクトはなんか微妙に意味が違ったりするので注意しないといけない．
ビコンパクトとかも．

論文
\cite{MR1512217}
は多分読めた．
可分距離空間がヒルベルトキューブに位相的に埋め込めることを示していると思う．


論文
\cite{MR1544691}
は読めなかった．
平面の連結集合の話をしている気がする．


論文
\cite{MR1512215}
はコンパクト空間が距離化可能であるための必要十分条件は
第二可算ということを証明していると思う．

論文
\cite{MR1512258}
は連結ハウスドルフ可算空間を構成していると思う．



論文
\cite{MR1512260}
はウリゾーンの距離化可能定理：
第二可算な正規空間は距離化可能である，
を証明している．
現代と変わらない証明が展開されてて
面白い．

論文
\cite{MR1544734}
は読めなかった．
この論文は普通に解析学の研究をしている．
ウリゾーンの一番最初の研究領域は解析学という話を聞いたことがあるような気がする．

論文
\cite{MR1512406}
は連続体の論文でconnected im kleinen (cik)
を使って何かしているようだ．
よめなかった．


%READ!
%myplain is the same to amsplain style. 
\nocite{*}
\bibliographystyle{myplaindoidoi}
\bibliography{Urysohn.bib}



\end{document}